\scanpage{44}
% Pencil archive foliation: 43.
par
\[
\operatorname{cl}(X)\unsure{\rightsquigarrow|\longmapsto}
\sum(-1)^i\operatorname{cl}H_!^i(X,\unsure{\mathbb{Q}|\mathbb{Z}}(0))
\]
où bien entendu $H_!^i(X,\mathbb{Z}(0))$ est
\deleted{considéré}\inserted{défini} comme motif sur $K$ par
$R^if_!(\unsure{\mathbb{Q}|\mathbb{Z}}(0))$ ;
$f:X\unsure{-|\longrightarrow}\operatorname{Spec}K$ la projection.
Prenant les composantes homogènes de
$\sigma\operatorname{cl}(X)$, on a des $\sigma_i(\operatorname{cl}(X))$ qui sont
des ``motifs virtuels'', et permettent de
définir des ``nbs de Betti virtuels'', ceux
demandés par Serre.

\deleted{
Si on filtre $L(K)$ par la dimension
en des $L_i(K)$, on trouve
\[
L_i(K)\longrightarrow D_i(\mathrm{M}^+_{\unsure{\ell|{}}}(K))
\]
On peut se demander si c'est surjectif,
ou injectif, \deleted{bijectif serait étonnant}
\unsure{même} mod torsion. Déjà pour $i=1$
($K$ alg clos) \inserted{\unsure{si le} degré de trans. infini du corps de base},
rien n'est connu.
On peut poser la question en
considérant des $X/K$ munis d'un groupe fini
donné $G$ d'autom. --- mais il faut
}
\margin{\deleted{
Pb pas bien \unsure{posé} car \unsure{aucune}
\unsure{définition} \uline{\unsure{rigide|rigidité}} \unsure{à|et}
telle quelle. cf plus bas la formulation
correcte \ldots
En fait
\[
L_i(K)\longrightarrow
\mathrm{M}_{(i)}^+(K)+\xi\mathrm{M}_{(i-1)}^+(K)
+\xi^2\mathrm{M}_{(i-2)}^+(K)+\cdots
\]
}}
